\documentclass[aspectratio=169,11pt]{beamer}
% Lecture 05 — Character data type
% Course: Introduction to Programming with Kotlin

% Shared Beamer preamble for Introduction to Programming with Kotlin
% Include with: % Shared Beamer preamble for Introduction to Programming with Kotlin
% Include with: % Shared Beamer preamble for Introduction to Programming with Kotlin
% Include with: \input{preamble}

\usetheme{Madrid}
\usecolortheme{default}

% Clean, professional palette (deep navy + muted teal accent)
\definecolor{LectureNavy}{HTML}{1B3A4B}
\definecolor{LectureAccent}{HTML}{2A9D8F}
\definecolor{LectureGray}{HTML}{4A5568}
\definecolor{CodeBg}{HTML}{F7FAFC}
\definecolor{AlertBg}{HTML}{FFF5F5}
\definecolor{TryBg}{HTML}{F0FFF4}
\definecolor{QuizBg}{HTML}{EBF8FF}
\definecolor{AnswerKeyBg}{HTML}{FFFFF0}
\definecolor{AnswerGold}{HTML}{B7791F}
\definecolor{KeywordBlue}{HTML}{2B6CB0}
\definecolor{StringGreen}{HTML}{276749}
\definecolor{CommentGray}{HTML}{718096}

\setbeamercolor{structure}{fg=LectureNavy}
\setbeamercolor{palette primary}{bg=LectureNavy,fg=white}
\setbeamercolor{palette secondary}{bg=LectureAccent,fg=white}
\setbeamercolor{palette tertiary}{bg=LectureGray,fg=white}
% Madrid puts title/subtitle in a coloured box (inherits palette primary bg).
% Title text must be light on that dark box — not navy-on-navy.
\setbeamercolor{title}{fg=white,bg=LectureNavy}
\setbeamercolor{subtitle}{fg=white,bg=LectureNavy}
\setbeamercolor{author}{fg=LectureNavy}
\setbeamercolor{date}{fg=LectureGray}
\setbeamercolor{institute}{fg=LectureGray}
\setbeamercolor{frametitle}{bg=LectureNavy,fg=white}
\setbeamercolor{block title}{bg=LectureNavy,fg=white}
\setbeamercolor{block body}{bg=CodeBg,fg=black}
\setbeamercolor{block title alerted}{bg=LectureAccent,fg=white}
\setbeamercolor{block body alerted}{bg=TryBg,fg=black}
\setbeamercolor{block title example}{bg=LectureGray,fg=white}
\setbeamercolor{block body example}{bg=AlertBg,fg=black}

\setbeamertemplate{navigation symbols}{}
\setbeamertemplate{itemize items}[circle]
\setbeamertemplate{enumerate items}[default]

\usepackage[T1]{fontenc}
\usepackage[utf8]{inputenc}
\usepackage{lmodern}
\usepackage{booktabs}
\usepackage{tabularx}
\usepackage{graphicx}
\usepackage{hyperref}
\usepackage{listings}
\usepackage{xcolor}
\usepackage{tikz}
\usetikzlibrary{positioning,arrows.meta}

% listings: Kotlin without shell-escape (no built-in Kotlin language)
\lstdefinelanguage{Kotlin}{
  morekeywords={
    abstract,actual,annotation,as,break,by,catch,class,companion,const,
    constructor,continue,crossinline,data,do,dynamic,else,enum,expect,
    external,false,field,file,final,finally,for,fun,get,if,import,in,
    infix,init,inline,inner,interface,internal,is,lateinit,noinline,null,
    object,open,operator,out,override,package,private,protected,public,
    reified,return,sealed,set,super,suspend,tailrec,this,throw,true,try,
    typealias,typeof,val,var,when,where,while
  },
  sensitive=true,
  morecomment=[l]{//},
  morecomment=[s]{/*}{*/},
  morestring=[b]",
  morestring=[b]',
  morestring=[s]{"""}{"""},
}

\lstdefinestyle{kotlinstyle}{
  language=Kotlin,
  basicstyle=\ttfamily\small,
  keywordstyle=\color{KeywordBlue}\bfseries,
  commentstyle=\color{CommentGray}\itshape,
  stringstyle=\color{StringGreen},
  numbers=left,
  numberstyle=\tiny\color{CommentGray},
  numbersep=6pt,
  frame=single,
  rulecolor=\color{LectureGray!40},
  backgroundcolor=\color{CodeBg},
  breaklines=true,
  showstringspaces=false,
  tabsize=2,
  captionpos=b,
  morekeywords={readln,readLine,println,print,listOf,mutableListOf,toInt,toDouble,toLong,toShort,toByte,toChar,toString,split,size},
}
\lstset{style=kotlinstyle}

\newcommand{\kotlincode}[1]{\texttt{\small #1}}
\newcommand{\trythis}[1]{%
  \begin{alertblock}{Try this}
    #1
  \end{alertblock}%
}
\newcommand{\commonmistake}[1]{%
  \begin{exampleblock}{Common mistake}
    #1
  \end{exampleblock}%
}
% Formative in-class checks: question slide, then a dedicated Answer slide
\newcommand{\quizpause}{%
  \par\vspace{0.5em}
  {\small\textit{Pause --- answer (clicker / raise hand / neighbour) before the next slide.}}%
}
\newcommand{\quizbox}[1]{%
  \setbeamercolor{block title}{bg=KeywordBlue,fg=white}%
  \setbeamercolor{block body}{bg=QuizBg,fg=black}%
  \begin{block}{Quiz}
    #1
  \end{block}%
  \setbeamercolor{block title}{bg=LectureNavy,fg=white}%
  \setbeamercolor{block body}{bg=CodeBg,fg=black}%
}
\newcommand{\answerbox}[1]{%
  \setbeamercolor{block title}{bg=AnswerGold,fg=white}%
  \setbeamercolor{block body}{bg=AnswerKeyBg,fg=black}%
  \begin{block}{Answer}
    #1
  \end{block}%
  \setbeamercolor{block title}{bg=LectureNavy,fg=white}%
  \setbeamercolor{block body}{bg=CodeBg,fg=black}%
}

% Empty author/date placeholders for the lecturer to fill
\author{}
\date{}


\usetheme{Madrid}
\usecolortheme{default}

% Clean, professional palette (deep navy + muted teal accent)
\definecolor{LectureNavy}{HTML}{1B3A4B}
\definecolor{LectureAccent}{HTML}{2A9D8F}
\definecolor{LectureGray}{HTML}{4A5568}
\definecolor{CodeBg}{HTML}{F7FAFC}
\definecolor{AlertBg}{HTML}{FFF5F5}
\definecolor{TryBg}{HTML}{F0FFF4}
\definecolor{QuizBg}{HTML}{EBF8FF}
\definecolor{AnswerKeyBg}{HTML}{FFFFF0}
\definecolor{AnswerGold}{HTML}{B7791F}
\definecolor{KeywordBlue}{HTML}{2B6CB0}
\definecolor{StringGreen}{HTML}{276749}
\definecolor{CommentGray}{HTML}{718096}

\setbeamercolor{structure}{fg=LectureNavy}
\setbeamercolor{palette primary}{bg=LectureNavy,fg=white}
\setbeamercolor{palette secondary}{bg=LectureAccent,fg=white}
\setbeamercolor{palette tertiary}{bg=LectureGray,fg=white}
% Madrid puts title/subtitle in a coloured box (inherits palette primary bg).
% Title text must be light on that dark box — not navy-on-navy.
\setbeamercolor{title}{fg=white,bg=LectureNavy}
\setbeamercolor{subtitle}{fg=white,bg=LectureNavy}
\setbeamercolor{author}{fg=LectureNavy}
\setbeamercolor{date}{fg=LectureGray}
\setbeamercolor{institute}{fg=LectureGray}
\setbeamercolor{frametitle}{bg=LectureNavy,fg=white}
\setbeamercolor{block title}{bg=LectureNavy,fg=white}
\setbeamercolor{block body}{bg=CodeBg,fg=black}
\setbeamercolor{block title alerted}{bg=LectureAccent,fg=white}
\setbeamercolor{block body alerted}{bg=TryBg,fg=black}
\setbeamercolor{block title example}{bg=LectureGray,fg=white}
\setbeamercolor{block body example}{bg=AlertBg,fg=black}

\setbeamertemplate{navigation symbols}{}
\setbeamertemplate{itemize items}[circle]
\setbeamertemplate{enumerate items}[default]

\usepackage[T1]{fontenc}
\usepackage[utf8]{inputenc}
\usepackage{lmodern}
\usepackage{booktabs}
\usepackage{tabularx}
\usepackage{graphicx}
\usepackage{hyperref}
\usepackage{listings}
\usepackage{xcolor}
\usepackage{tikz}
\usetikzlibrary{positioning,arrows.meta}

% listings: Kotlin without shell-escape (no built-in Kotlin language)
\lstdefinelanguage{Kotlin}{
  morekeywords={
    abstract,actual,annotation,as,break,by,catch,class,companion,const,
    constructor,continue,crossinline,data,do,dynamic,else,enum,expect,
    external,false,field,file,final,finally,for,fun,get,if,import,in,
    infix,init,inline,inner,interface,internal,is,lateinit,noinline,null,
    object,open,operator,out,override,package,private,protected,public,
    reified,return,sealed,set,super,suspend,tailrec,this,throw,true,try,
    typealias,typeof,val,var,when,where,while
  },
  sensitive=true,
  morecomment=[l]{//},
  morecomment=[s]{/*}{*/},
  morestring=[b]",
  morestring=[b]',
  morestring=[s]{"""}{"""},
}

\lstdefinestyle{kotlinstyle}{
  language=Kotlin,
  basicstyle=\ttfamily\small,
  keywordstyle=\color{KeywordBlue}\bfseries,
  commentstyle=\color{CommentGray}\itshape,
  stringstyle=\color{StringGreen},
  numbers=left,
  numberstyle=\tiny\color{CommentGray},
  numbersep=6pt,
  frame=single,
  rulecolor=\color{LectureGray!40},
  backgroundcolor=\color{CodeBg},
  breaklines=true,
  showstringspaces=false,
  tabsize=2,
  captionpos=b,
  morekeywords={readln,readLine,println,print,listOf,mutableListOf,toInt,toDouble,toLong,toShort,toByte,toChar,toString,split,size},
}
\lstset{style=kotlinstyle}

\newcommand{\kotlincode}[1]{\texttt{\small #1}}
\newcommand{\trythis}[1]{%
  \begin{alertblock}{Try this}
    #1
  \end{alertblock}%
}
\newcommand{\commonmistake}[1]{%
  \begin{exampleblock}{Common mistake}
    #1
  \end{exampleblock}%
}
% Formative in-class checks: question slide, then a dedicated Answer slide
\newcommand{\quizpause}{%
  \par\vspace{0.5em}
  {\small\textit{Pause --- answer (clicker / raise hand / neighbour) before the next slide.}}%
}
\newcommand{\quizbox}[1]{%
  \setbeamercolor{block title}{bg=KeywordBlue,fg=white}%
  \setbeamercolor{block body}{bg=QuizBg,fg=black}%
  \begin{block}{Quiz}
    #1
  \end{block}%
  \setbeamercolor{block title}{bg=LectureNavy,fg=white}%
  \setbeamercolor{block body}{bg=CodeBg,fg=black}%
}
\newcommand{\answerbox}[1]{%
  \setbeamercolor{block title}{bg=AnswerGold,fg=white}%
  \setbeamercolor{block body}{bg=AnswerKeyBg,fg=black}%
  \begin{block}{Answer}
    #1
  \end{block}%
  \setbeamercolor{block title}{bg=LectureNavy,fg=white}%
  \setbeamercolor{block body}{bg=CodeBg,fg=black}%
}

% Empty author/date placeholders for the lecturer to fill
\author{}
\date{}


\usetheme{Madrid}
\usecolortheme{default}

% Clean, professional palette (deep navy + muted teal accent)
\definecolor{LectureNavy}{HTML}{1B3A4B}
\definecolor{LectureAccent}{HTML}{2A9D8F}
\definecolor{LectureGray}{HTML}{4A5568}
\definecolor{CodeBg}{HTML}{F7FAFC}
\definecolor{AlertBg}{HTML}{FFF5F5}
\definecolor{TryBg}{HTML}{F0FFF4}
\definecolor{QuizBg}{HTML}{EBF8FF}
\definecolor{AnswerKeyBg}{HTML}{FFFFF0}
\definecolor{AnswerGold}{HTML}{B7791F}
\definecolor{KeywordBlue}{HTML}{2B6CB0}
\definecolor{StringGreen}{HTML}{276749}
\definecolor{CommentGray}{HTML}{718096}

\setbeamercolor{structure}{fg=LectureNavy}
\setbeamercolor{palette primary}{bg=LectureNavy,fg=white}
\setbeamercolor{palette secondary}{bg=LectureAccent,fg=white}
\setbeamercolor{palette tertiary}{bg=LectureGray,fg=white}
% Madrid puts title/subtitle in a coloured box (inherits palette primary bg).
% Title text must be light on that dark box — not navy-on-navy.
\setbeamercolor{title}{fg=white,bg=LectureNavy}
\setbeamercolor{subtitle}{fg=white,bg=LectureNavy}
\setbeamercolor{author}{fg=LectureNavy}
\setbeamercolor{date}{fg=LectureGray}
\setbeamercolor{institute}{fg=LectureGray}
\setbeamercolor{frametitle}{bg=LectureNavy,fg=white}
\setbeamercolor{block title}{bg=LectureNavy,fg=white}
\setbeamercolor{block body}{bg=CodeBg,fg=black}
\setbeamercolor{block title alerted}{bg=LectureAccent,fg=white}
\setbeamercolor{block body alerted}{bg=TryBg,fg=black}
\setbeamercolor{block title example}{bg=LectureGray,fg=white}
\setbeamercolor{block body example}{bg=AlertBg,fg=black}

\setbeamertemplate{navigation symbols}{}
\setbeamertemplate{itemize items}[circle]
\setbeamertemplate{enumerate items}[default]

\usepackage[T1]{fontenc}
\usepackage[utf8]{inputenc}
\usepackage{lmodern}
\usepackage{booktabs}
\usepackage{tabularx}
\usepackage{graphicx}
\usepackage{hyperref}
\usepackage{listings}
\usepackage{xcolor}
\usepackage{tikz}
\usetikzlibrary{positioning,arrows.meta}

% listings: Kotlin without shell-escape (no built-in Kotlin language)
\lstdefinelanguage{Kotlin}{
  morekeywords={
    abstract,actual,annotation,as,break,by,catch,class,companion,const,
    constructor,continue,crossinline,data,do,dynamic,else,enum,expect,
    external,false,field,file,final,finally,for,fun,get,if,import,in,
    infix,init,inline,inner,interface,internal,is,lateinit,noinline,null,
    object,open,operator,out,override,package,private,protected,public,
    reified,return,sealed,set,super,suspend,tailrec,this,throw,true,try,
    typealias,typeof,val,var,when,where,while
  },
  sensitive=true,
  morecomment=[l]{//},
  morecomment=[s]{/*}{*/},
  morestring=[b]",
  morestring=[b]',
  morestring=[s]{"""}{"""},
}

\lstdefinestyle{kotlinstyle}{
  language=Kotlin,
  basicstyle=\ttfamily\small,
  keywordstyle=\color{KeywordBlue}\bfseries,
  commentstyle=\color{CommentGray}\itshape,
  stringstyle=\color{StringGreen},
  numbers=left,
  numberstyle=\tiny\color{CommentGray},
  numbersep=6pt,
  frame=single,
  rulecolor=\color{LectureGray!40},
  backgroundcolor=\color{CodeBg},
  breaklines=true,
  showstringspaces=false,
  tabsize=2,
  captionpos=b,
  morekeywords={readln,readLine,println,print,listOf,mutableListOf,toInt,toDouble,toLong,toShort,toByte,toChar,toString,split,size},
}
\lstset{style=kotlinstyle}

\newcommand{\kotlincode}[1]{\texttt{\small #1}}
\newcommand{\trythis}[1]{%
  \begin{alertblock}{Try this}
    #1
  \end{alertblock}%
}
\newcommand{\commonmistake}[1]{%
  \begin{exampleblock}{Common mistake}
    #1
  \end{exampleblock}%
}
% Formative in-class checks: question slide, then a dedicated Answer slide
\newcommand{\quizpause}{%
  \par\vspace{0.5em}
  {\small\textit{Pause --- answer (clicker / raise hand / neighbour) before the next slide.}}%
}
\newcommand{\quizbox}[1]{%
  \setbeamercolor{block title}{bg=KeywordBlue,fg=white}%
  \setbeamercolor{block body}{bg=QuizBg,fg=black}%
  \begin{block}{Quiz}
    #1
  \end{block}%
  \setbeamercolor{block title}{bg=LectureNavy,fg=white}%
  \setbeamercolor{block body}{bg=CodeBg,fg=black}%
}
\newcommand{\answerbox}[1]{%
  \setbeamercolor{block title}{bg=AnswerGold,fg=white}%
  \setbeamercolor{block body}{bg=AnswerKeyBg,fg=black}%
  \begin{block}{Answer}
    #1
  \end{block}%
  \setbeamercolor{block title}{bg=LectureNavy,fg=white}%
  \setbeamercolor{block body}{bg=CodeBg,fg=black}%
}

% Empty author/date placeholders for the lecturer to fill
\author{}
\date{}


\title{Lecture 05: The Character Type}
\subtitle{Introduction to Programming with Kotlin}

\begin{document}

%----------------------------------------------------------------------
    \begin{frame}
        \titlepage
    \end{frame}

%----------------------------------------------------------------------
    \begin{frame}{Agenda}
        \begin{enumerate}
            \item What is \kotlincode{Char}?
            \item Characters vs strings
            \item Codes, ranges, and classification
            \item Working with characters in text
        \end{enumerate}
    \end{frame}

%======================================================================
    \section{Character data type}
%======================================================================

    \begin{frame}{The \kotlincode{Char} type}
        A \kotlincode{Char} holds a \textbf{single} Unicode character.
        \vspace{0.6em}
        \begin{itemize}
            \item Literals use single quotes: \kotlincode{'A'}, \kotlincode{'7'}, \kotlincode{'*'}
            \item Strings use double quotes: \kotlincode{"A"} is a \kotlincode{String} of length 1
            \item Useful for letters, digits, punctuation, and simple text processing
        \end{itemize}
    \end{frame}

    \begin{frame}[fragile]{Declaring characters}
        \begin{lstlisting}
fun main() {
    val grade: Char = 'B'
    val digit = '7'
    val newline = '\n'
    val quote = '\''

    println(grade)
    println("Digit char: $digit")
}
        \end{lstlisting}
        \vspace{0.3em}
        Escapes work inside character literals too (\kotlincode{'\textbackslash n'},
        \kotlincode{'\textbackslash t'}, \kotlincode{'\textbackslash''}).
    \end{frame}

    \begin{frame}{Char vs String}
        \begin{tabular}{lll}
            \toprule
            & \kotlincode{Char} & \kotlincode{String} \\
            \midrule
            Quotes & \kotlincode{'A'} & \kotlincode{"A"} / \kotlincode{"Hi"} \\
            Length & always one char & zero or more \\
            Indexing & --- & \kotlincode{s[i]} yields \kotlincode{Char} \\
            Typical use & symbols, letters & words, lines, messages \\
            \bottomrule
        \end{tabular}
        \vspace{0.6em}
        \commonmistake{Writing \kotlincode{"A"} when an API expects \kotlincode{Char},
        or the reverse. The types are different.}
    \end{frame}

    \begin{frame}[fragile]{Characters from strings}
        \begin{lstlisting}
fun main() {
    val word = readln()
    if (word.isNotEmpty()) {
        val first: Char = word[0]
        val last = word[word.length - 1]
        println("First: $first, last: $last")
    }
}
        \end{lstlisting}
        \vspace{0.3em}
        Strings also have \kotlincode{.length} (like \kotlincode{.size} on lists).
        \kotlincode{word.first()} / \kotlincode{word.last()} are handy when non-empty.
    \end{frame}

    %--- Quiz checkpoint 1 ---
    \begin{frame}{Quiz 1 --- Char basics}
        \quizbox{Which literal has type \kotlincode{Char}?}
        \begin{enumerate}
            \item \kotlincode{"Z"}
            \item \kotlincode{'Z'}
            \item \kotlincode{Z}
            \item \kotlincode{char("Z")}
        \end{enumerate}
        \quizpause
    \end{frame}

    \begin{frame}{Answer --- Quiz 1}
        \answerbox{\textbf{2.} \kotlincode{'Z'}}
        \vspace{0.4em}
        Single quotes create a \kotlincode{Char}.
        Double quotes create a \kotlincode{String}.
    \end{frame}

%======================================================================
    \section{Codes and classification}
%======================================================================

    \begin{frame}{Unicode code points (intuition)}
        Each character has an integer code.
        \vspace{0.5em}
        \begin{itemize}
            \item Digits \kotlincode{'0'}..\kotlincode{'9'} are consecutive
            \item Uppercase \kotlincode{'A'}..\kotlincode{'Z'} are consecutive
            \item Lowercase \kotlincode{'a'}..\kotlincode{'z'} are consecutive
        \end{itemize}
        \vspace{0.4em}
        You can convert: \kotlincode{c.code} (Int) and \kotlincode{Char(code)}.
    \end{frame}

    \begin{frame}[fragile]{Codes and arithmetic}
        \begin{lstlisting}
fun main() {
    val c = 'A'
    println(c.code)           // 65
    println(Char(66))         // B

    val next = c + 1          // Char arithmetic -> 'B'
    println(next)

    val digit = '7'
    val value = digit - '0'   // 7 as Int
    println(value)
}
        \end{lstlisting}
        \kotlincode{digit - '0'} is a classic way to turn a digit character into its
        numeric value $0$--$9$.
    \end{frame}

    \begin{frame}[fragile]{Ranges of characters}
        \begin{lstlisting}
fun main() {
    val ch = readln()[0]
    when (ch) {
        in '0'..'9' -> println("Digit")
        in 'A'..'Z' -> println("Uppercase")
        in 'a'..'z' -> println("Lowercase")
        else -> println("Other")
    }
}
        \end{lstlisting}
        \trythis{Count how many letters are in an input line using a \kotlincode{for} loop
        over \kotlincode{line} (strings are iterable as characters).}
    \end{frame}

    \begin{frame}[fragile]{Built-in checks}
        Kotlin provides helpers on \kotlincode{Char}:
        \vspace{0.3em}
        \begin{lstlisting}
fun main() {
    val ch = readln()[0]
    println(ch.isDigit())
    println(ch.isLetter())
    println(ch.isLetterOrDigit())
    println(ch.isWhitespace())
    println(ch.lowercaseChar())
    println(ch.uppercaseChar())
}
        \end{lstlisting}
        Prefer these for clarity when they match your need.
    \end{frame}

    %--- Quiz checkpoint 2 ---
    \begin{frame}{Quiz 2 --- Digit value}
        \quizbox{What is \kotlincode{'5' - '0'} in Kotlin?}
        \begin{enumerate}
            \item \kotlincode{'5'}
            \item \kotlincode{50}
            \item \kotlincode{5}
            \item A compile error
        \end{enumerate}
        \quizpause
    \end{frame}

    \begin{frame}{Answer --- Quiz 2}
        \answerbox{\textbf{3.} \kotlincode{5}}
        \vspace{0.4em}
        Subtracting two \kotlincode{Char} values yields the distance between their codes
        as an \kotlincode{Int}. For digit characters, that distance is the digit's value.
    \end{frame}

%======================================================================
    \section{Characters in practice}
%======================================================================

    \begin{frame}[fragile]{Iterate characters of a string}
        \begin{lstlisting}
fun main() {
    val line = readln()
    var letters = 0
    for (ch in line) {
        if (ch.isLetter()) {
            letters += 1
        }
    }
    println("Letters: $letters")
}
        \end{lstlisting}
    \end{frame}

    \begin{frame}[fragile]{Build a string from characters}
        \begin{lstlisting}
fun main() {
    val word = readln()
    var onlyLetters = ""
    for (ch in word) {
        if (ch.isLetter()) {
            onlyLetters += ch
        }
    }
    println(onlyLetters)
}
        \end{lstlisting}
        \vspace{0.3em}
        For heavy concatenation, later you will meet \kotlincode{StringBuilder};
        for intro exercises, \kotlincode{+=} is fine.
    \end{frame}

    \begin{frame}[fragile]{Caesar-style shift (tiny demo)}
        \begin{lstlisting}
fun main() {
    val ch = readln()[0]
    if (ch in 'a'..'z') {
        val shifted = 'a' + (ch - 'a' + 1) % 26
        println(shifted)
    } else {
        println(ch)
    }
}
        \end{lstlisting}
        Shifts \kotlincode{'a'}$\rightarrow$\kotlincode{'b'}, \ldots,
        \kotlincode{'z'}$\rightarrow$\kotlincode{'a'}.
    \end{frame}

    %--- Quiz checkpoint 3 ---
    \begin{frame}{Quiz 3 --- Iteration}
        \quizbox{After this code, what is \kotlincode{n}?\\
        \kotlincode{var n = 0; for (ch in "A1b") if (ch.isLetter()) n++}}
        \begin{enumerate}
            \item 0
            \item 1
            \item 2
            \item 3
        \end{enumerate}
        \quizpause
    \end{frame}

    \begin{frame}{Answer --- Quiz 3}
        \answerbox{\textbf{3.} 2}
        \vspace{0.4em}
        \kotlincode{'A'} and \kotlincode{'b'} are letters; \kotlincode{'1'} is not.
    \end{frame}

%======================================================================
    \section{Wrap-up}
%======================================================================

    \begin{frame}{Summary}
        \begin{itemize}
            \item \kotlincode{Char} --- one character in single quotes
            \item Strings are sequences of characters; \kotlincode{s[i]} is a \kotlincode{Char}
            \item Codes and ranges enable digit/letter checks and simple ciphers
            \item Prefer \kotlincode{isDigit()}, \kotlincode{isLetter()}, \ldots when appropriate
        \end{itemize}
    \end{frame}

    \begin{frame}{Course checkpoint}
        You can now write Kotlin programs that:
        \begin{itemize}
            \item Read input and use variables / templates
            \item Compute with integers and floats
            \item Branch with \kotlincode{if} / \kotlincode{when}
            \item Loop over lists, ranges, and characters
            \item Split lines and process text at character level
        \end{itemize}
        Next steps in a fuller course: functions, null-safety, classes, and testing.
    \end{frame}

\end{document}
